% =====================================================================
% References -- consolidated from every chapter's Further Reading section.
% GENERATED by a script from a curated list; edit the list, not by hand.
% =====================================================================
\chapter*{References}
\addcontentsline{toc}{chapter}{References}
\markboth{References}{References}
\label{references}

\providecommand{\citedin}[1]{{\unskip\nobreak\hskip1.5em plus 1fil\penalty50
  \hskip0pt plus -1fil\hbox{}\nobreak\hskip0pt plus 1fil
  \mbox{\footnotesize\itshape #1}\parfillskip=0pt\finalhyphendemerits=0\par}}

Every chapter of this book ends with a \emph{Further Reading} section in which each work is
annotated: what it covers, why it is worth your time, and which part of the chapter it
supports. This list collects every work cited in those sections into one alphabetized
reference, so that you can find a book or paper again without remembering which chapter
recommended it. At the right of each entry are the chapters (and appendices) that cite it;
the numbers are links, and the annotations are to be found in those chapters. Books, papers,
standards and other sources are listed separately; standards are ordered by organization and
then by document number. Standards bodies revise their documents continually, so unless an
edition is stated, consult the current version.

\section*{Books}
\addcontentsline{toc}{section}{Books}
\markright{Books}
\begingroup\small
\begin{itemize}[leftmargin=1.5em,itemindent=-1.5em,itemsep=1.5pt plus 1pt,parsep=0pt,topsep=4pt,label={}]\rightskip=0pt plus 3em\relax
\item M.~Abramowitz and I.~A.~Stegun (eds.), \emph{Handbook of Mathematical Functions}, 1964.\citedin{App.~\hyperref[appA]{A}}
\item G.~P.~Agrawal, \emph{Fiber-Optic Communication Systems}, 4th ed., Wiley, 2010 (5th ed., 2021).\citedin{Ch.~\hyperref[ch:wireline]{24}}
\item G.~P.~Agrawal, \emph{Nonlinear Fiber Optics}.\citedin{Ch.~\hyperref[ch:wireline]{24}}
\item J.~B.~Anderson, T.~Aulin and C.-E.~Sundberg, \emph{Digital Phase Modulation}, Plenum, 1986.\citedin{Ch.~\hyperref[ch:modulation]{9}}
\item American Radio Relay League, \emph{The ARRL Handbook for Radio Communications}, ARRL (annual).\citedin{Ch.~\hyperref[ch:analog]{4}}
\item C.~A.~Balanis, \emph{Antenna Theory: Analysis and Design}, 4th ed., Wiley, 2016.\citedin{Ch.~\hyperref[ch:channels]{11}}
\item J.~R.~Barry, E.~A.~Lee and D.~G.~Messerschmitt, \emph{Digital Communication}, 3rd ed., Kluwer, 2004 (2nd ed., by E.~A.~Lee and D.~G.~Messerschmitt, Kluwer, 1994).\citedin{Chs.~\hyperref[ch:baseband]{8}, \hyperref[ch:equalization]{12}}
\item Members of the Technical Staff, Bell Telephone Laboratories, \emph{Transmission Systems for Communications}, 5th ed., 1982.\citedin{Ch.~\hyperref[ch:analog]{4}}
\item J.~Bellamy, \emph{Digital Telephony}, 3rd ed., Wiley, 2000.\citedin{Chs.~\hyperref[ch:pcm]{5}, \hyperref[ch:baseband]{8}, \hyperref[ch:wireline]{24}}
\item K.~B.~Benson (ed.), \emph{Television Engineering Handbook}, McGraw-Hill, 1986.\citedin{Ch.~\hyperref[ch:analog]{4}}
\item D.~Bertsekas and R.~Gallager, \emph{Data Networks}, 2nd ed., Prentice Hall, 1992.\citedin{Ch.~\hyperref[ch:multipleaccess]{20}}
\item J.~W.~Betz, \emph{Engineering Satellite-Based Navigation and Timing: Global Navigation Satellite Systems, Signals, and Receivers}, Wiley-IEEE Press, 2016.\citedin{Ch.~\hyperref[ch:spreadspectrum]{18}}
\item E.~Bj\"ornson, J.~Hoydis and L.~Sanguinetti, \emph{Massive MIMO Networks: Spectral, Energy, and Hardware Efficiency}, Foundations and Trends in Signal Processing, 2017.\citedin{Ch.~\hyperref[ch:mimo]{19}}
\item R.~E.~Blahut, \emph{Algebraic Codes for Data Transmission}, Cambridge University Press, 2003.\citedin{Ch.~\hyperref[ch:classiccodes]{14}}
\item K.~Borre, D.~M.~Akos, N.~Bertelsen, P.~Rinder and S.~H.~Jensen, \emph{A Software-Defined GPS and Galileo Receiver: A Single-Frequency Approach}, Birkh\"auser, 2007.\citedin{Ch.~\hyperref[ch:spreadspectrum]{18}}
\item S.~Boyd and L.~Vandenberghe, \emph{Convex Optimization}, Cambridge University Press, 2004.\citedin{App.~\hyperref[appA]{A}}
\item R.~N.~Bracewell, \emph{The Fourier Transform and Its Applications}, 3rd ed., McGraw-Hill, 2000.\citedin{Ch.~\hyperref[ch:signals]{2}; App.~\hyperref[appA]{A}}
\item T.~F.~Collins, R.~Getz, D.~Pu and A.~M.~Wyglinski, \emph{Software-Defined Radio for Engineers}, Artech House, 2018.\citedin{App.~\hyperref[appB]{B}}
\item T.~M.~Cover and J.~A.~Thomas, \emph{Elements of Information Theory}, 2nd ed., Wiley, 2006.\citedin{Ch.~\hyperref[ch:infotheory]{13}}
\item S.~C.~Cripps, \emph{RF Power Amplifiers for Wireless Communications}, 2nd ed., Artech House, 2006.\citedin{Ch.~\hyperref[ch:sdr]{7}}
\item R.~E.~Crochiere and L.~R.~Rabiner, \emph{Multirate Digital Signal Processing}, Prentice Hall, 1983.\citedin{Ch.~\hyperref[ch:dsp]{6}}
\item E.~Dahlman, S.~Parkvall and J.~Sk\"old, \emph{4G LTE-Advanced Pro and the Road to 5G}, 3rd ed., Academic Press, 2016.\citedin{Ch.~\hyperref[ch:cellular]{21}}
\item E.~Dahlman, S.~Parkvall and J.~Sk\"old, \emph{5G NR: The Next Generation Wireless Access Technology}, 2nd ed., Academic Press, 2020.\citedin{Chs.~\hyperref[ch:ofdm]{17}, \hyperref[ch:multipleaccess]{20}, \hyperref[ch:cellular]{21}}
\item \"O.~T.~Demir, E.~Bj\"ornson and L.~Sanguinetti, \emph{Foundations of User-Centric Cell-Free Massive MIMO}, Foundations and Trends in Signal Processing, Now Publishers, 2021.\citedin{Ch.~\hyperref[ch:sixg]{25}}
\item D.~M.~Dobkin, \emph{The RF in RFID: UHF RFID in Practice}, 2nd ed., Newnes, 2012.\citedin{Ch.~\hyperref[ch:wifi]{22}}
\item A.~El~Gamal and Y.-H.~Kim, \emph{Network Information Theory}, Cambridge University Press, 2011.\citedin{Ch.~\hyperref[ch:infotheory]{13}}
\item K.~Finkenzeller, \emph{RFID Handbook}, 3rd ed., Wiley, 2010.\citedin{Ch.~\hyperref[ch:wifi]{22}}
\item R.~F.~H.~Fischer, \emph{Precoding and Signal Shaping for Digital Transmission}, Wiley, 2002.\citedin{Ch.~\hyperref[ch:equalization]{12}}
\item R.~G.~Gallager, \emph{Low-Density Parity-Check Codes}, MIT Press, 1963.\citedin{Ch.~\hyperref[ch:moderncodes]{15}}
\item R.~G.~Gallager, \emph{Information Theory and Reliable Communication}, Wiley, 1968.\citedin{Ch.~\hyperref[ch:infotheory]{13}}
\item W.~A.~Gardner, \emph{Statistical Spectral Analysis: A Nonprobabilistic Theory}, Prentice Hall, 1988.\citedin{Ch.~\hyperref[ch:noise]{3}}
\item F.~M.~Gardner, \emph{Phaselock Techniques}, 3rd ed., Wiley, 2005.\citedin{Ch.~\hyperref[ch:sync]{10}}
\item M.~S.~Gast, \emph{802.11 Wireless Networks: The Definitive Guide}, 2nd ed., O'Reilly, 2005.\citedin{Ch.~\hyperref[ch:wifi]{22}}
\item M.~S.~Gast, \emph{802.11ac: A Survival Guide}, O'Reilly, 2013.\citedin{Ch.~\hyperref[ch:wifi]{22}}
\item A.~Gersho and R.~M.~Gray, \emph{Vector Quantization and Signal Compression}, Kluwer, 1992.\citedin{Ch.~\hyperref[ch:sourcecoding]{16}}
\item J.~Gertner, \emph{The Idea Factory: Bell Labs and the Great Age of American Innovation}, Penguin, 2012.\citedin{Ch.~\hyperref[ch:history]{1}}
\item J.~Gleick, \emph{The Information: A History, a Theory, a Flood}, Pantheon, 2011.\citedin{Ch.~\hyperref[ch:history]{1}}
\item I.~S.~Gradshteyn and I.~M.~Ryzhik, \emph{Table of Integrals, Series, and Products}, 8th ed., Academic Press, 2014.\citedin{App.~\hyperref[appA]{A}}
\item M.~Haenggi, \emph{Stochastic Geometry for Wireless Networks}, Cambridge University Press, 2012.\citedin{Ch.~\hyperref[ch:multipleaccess]{20}}
\item K.~Hafner and M.~Lyon, \emph{Where Wizards Stay Up Late: The Origins of the Internet}, Simon \& Schuster, 1996.\citedin{Ch.~\hyperref[ch:history]{1}}
\item f.~j.~harris, \emph{Multirate Signal Processing for Communication Systems}, 2nd ed., River Publishers, 2021.\citedin{Ch.~\hyperref[ch:dsp]{6}}
\item S.~Haykin and B.~Van~Veen, \emph{Signals and Systems}, 2nd ed., Wiley, 2003.\citedin{Ch.~\hyperref[ch:signals]{2}}
\item S.~Haykin and M.~Moher, \emph{Communication Systems}, 5th ed., Wiley, 2009.\citedin{Chs.~\hyperref[ch:signals]{2}, \hyperref[ch:analog]{4}}
\item S.~Haykin, \emph{Adaptive Filter Theory}, 5th ed., Pearson, 2014.\citedin{Ch.~\hyperref[ch:equalization]{12}}
\item R.~W.~Heath Jr.\ and A.~Lozano, \emph{Foundations of MIMO Communication}, Cambridge University Press, 2018.\citedin{Ch.~\hyperref[ch:mimo]{19}}
\item F.~Hillebrand (ed.), \emph{GSM and UMTS: The Creation of Global Mobile Communication}, Wiley, 2002.\citedin{Ch.~\hyperref[ch:cellular]{21}}
\item A.~Hj{\o}rungnes, \emph{Complex-Valued Matrix Derivatives}, Cambridge University Press, 2011.\citedin{App.~\hyperref[appA]{A}}
\item H.~Holma and A.~Toskala (eds.), \emph{WCDMA for UMTS: HSPA Evolution and LTE}, 5th ed., Wiley, 2010.\citedin{Ch.~\hyperref[ch:cellular]{21}}
\item G.~J.~Holzmann and B.~Pehrson, \emph{The Early History of Data Networks}, IEEE Computer Society Press, 1995.\citedin{Ch.~\hyperref[ch:history]{1}}
\item S.~Hong, \emph{Wireless: From Marconi's Black-Box to the Audion}, MIT Press, 2001.\citedin{Ch.~\hyperref[ch:history]{1}}
\item R.~A.~Horn and C.~R.~Johnson, \emph{Matrix Analysis}, 2nd ed., Cambridge University Press, 2013.\citedin{App.~\hyperref[appA]{A}}
\item A.~A.~Huurdeman, \emph{The Worldwide History of Telecommunications}, Wiley, 2003.\citedin{Ch.~\hyperref[ch:history]{1}}
\item L.~J.~Ippolito, \emph{Satellite Communications Systems Engineering: Atmospheric Effects, Satellite Link Design and System Performance}, 2nd ed., Wiley, 2017.\citedin{Ch.~\hyperref[ch:satellite]{23}}
\item L.~B.~Jackson, \emph{Digital Filters and Signal Processing}, 3rd ed., Kluwer, 1996.\citedin{Ch.~\hyperref[ch:dsp]{6}}
\item W.~C.~Jakes (ed.), \emph{Microwave Mobile Communications}, Wiley, 1974; IEEE Press reissue, 1994.\citedin{Ch.~\hyperref[ch:channels]{11}}
\item N.~S.~Jayant and P.~Noll, \emph{Digital Coding of Waveforms: Principles and Applications to Speech and Video}, Prentice Hall, 1984.\citedin{Chs.~\hyperref[ch:pcm]{5}, \hyperref[ch:sourcecoding]{16}}
\item R.~Johannesson and K.~Sh.~Zigangirov, \emph{Fundamentals of Convolutional Coding}, IEEE Press, 1999.\citedin{Ch.~\hyperref[ch:classiccodes]{14}}
\item E.~D.~Kaplan and C.~J.~Hegarty (eds.), \emph{Understanding GPS/GNSS: Principles and Applications}, 3rd ed., Artech House, 2017.\citedin{Ch.~\hyperref[ch:spreadspectrum]{18}}
\item S.~M.~Kay, \emph{Fundamentals of Statistical Signal Processing}, vol.~I, \emph{Estimation Theory} (1993), and vol.~II, \emph{Detection Theory} (1998), Prentice Hall.\citedin{Ch.~\hyperref[ch:noise]{3}}
\item W.~Kester (ed.), \emph{Data Conversion Handbook}, Analog Devices and Newnes, 2005.\citedin{Ch.~\hyperref[ch:pcm]{5}}
\item L.~Kleinrock, \emph{Queueing Systems, Volume~1: Theory}, Wiley, 1975.\citedin{Ch.~\hyperref[ch:multipleaccess]{20}}
\item A.~M.~Kondoz, \emph{Digital Speech: Coding for Low Bit Rate Communication Systems}, 2nd ed., Wiley, 2004.\citedin{Ch.~\hyperref[ch:sourcecoding]{16}}
\item B.~P.~Lathi and R.~Green, \emph{Linear Systems and Signals}, 3rd ed., Oxford University Press, 2018.\citedin{Ch.~\hyperref[ch:signals]{2}}
\item B.~P.~Lathi and Z.~Ding, \emph{Modern Digital and Analog Communication Systems}, 5th ed., Oxford University Press, 2019.\citedin{Ch.~\hyperref[ch:analog]{4}}
\item A.~Leon-Garcia, \emph{Probability, Statistics, and Random Processes for Electrical Engineering}, 3rd ed., Pearson, 2008.\citedin{Ch.~\hyperref[ch:noise]{3}}
\item T.~Lewis, \emph{Empire of the Air: The Men Who Made Radio}, HarperCollins, 1991.\citedin{Ch.~\hyperref[ch:analog]{4}}
\item Y.~Li and G.~L.~St\"uber (eds.), \emph{Orthogonal Frequency Division Multiplexing for Wireless Communications}, Springer, 2006.\citedin{Ch.~\hyperref[ch:ofdm]{17}}
\item M.~P.~Li, \emph{Jitter, Noise, and Signal Integrity at High-Speed}, Prentice Hall, 2007.\citedin{Ch.~\hyperref[ch:baseband]{8}}
\item S.~Lin and D.~J.~Costello, Jr., \emph{Error Control Coding}, 2nd ed., Pearson Prentice Hall, 2004.\citedin{Ch.~\hyperref[ch:classiccodes]{14}}
\item R.~G.~Lyons, \emph{Understanding Digital Signal Processing}, 3rd ed., Prentice Hall/Pearson, 2010.\citedin{Chs.~\hyperref[ch:signals]{2}, \hyperref[ch:dsp]{6}}
\item D.~J.~C.~MacKay, \emph{Information Theory, Inference, and Learning Algorithms}, Cambridge University Press, 2003.\citedin{Ch.~\hyperref[ch:infotheory]{13}}
\item F.~J.~MacWilliams and N.~J.~A.~Sloane, \emph{The Theory of Error-Correcting Codes}, North-Holland, 1977.\citedin{Ch.~\hyperref[ch:classiccodes]{14}}
\item G.~Maral, M.~Bousquet and Z.~Sun, \emph{Satellite Communications Systems: Systems, Techniques and Technology}, 6th ed., Wiley, 2020.\citedin{Ch.~\hyperref[ch:satellite]{23}}
\item T.~L.~Marzetta, E.~G.~Larsson, H.~Yang and H.~Q.~Ngo, \emph{Fundamentals of Massive MIMO}, Cambridge University Press, 2016.\citedin{Ch.~\hyperref[ch:mimo]{19}}
\item U.~Mengali and A.~N.~D'Andrea, \emph{Synchronization Techniques for Digital Receivers}, Plenum, 1997.\citedin{Ch.~\hyperref[ch:sync]{10}}
\item H.~Meyr, M.~Moeneclaey and S.~A.~Fechtel, \emph{Digital Communication Receivers: Synchronization, Channel Estimation, and Signal Processing}, Wiley, 1998.\citedin{Ch.~\hyperref[ch:sync]{10}}
\item P.~Misra and P.~Enge, \emph{Global Positioning System: Signals, Measurements, and Performance}, 2nd ed., Ganga-Jamuna Press, 2006.\citedin{Ch.~\hyperref[ch:spreadspectrum]{18}}
\item A.~F.~Molisch, \emph{Wireless Communications}, 2nd ed., Wiley--IEEE, 2011.\citedin{Ch.~\hyperref[ch:channels]{11}}
\item T.~K.~Moon, \emph{Error Correction Coding: Mathematical Methods and Algorithms}, Wiley, 2005.\citedin{Ch.~\hyperref[ch:classiccodes]{14}}
\item M.~Mouly and M.-B.~Pautet, \emph{The GSM System for Mobile Communications}, published by the authors, 1992.\citedin{Ch.~\hyperref[ch:cellular]{21}}
\item P.~J.~Nahin, \emph{Oliver Heaviside: The Life, Work, and Times of an Electrical Genius of the Victorian Age}, Johns Hopkins, 2002 (first published 1988).\citedin{Ch.~\hyperref[ch:history]{1}}
\item F.~W.~J.~Olver et al.\ (eds.), \emph{NIST Digital Library of Mathematical Functions}, \url{https://dlmf.nist.gov}.\citedin{App.~\hyperref[appA]{A}}
\item A.~V.~Oppenheim and A.~S.~Willsky, with S.~H.~Nawab, \emph{Signals and Systems}, 2nd ed., Prentice Hall, 1997.\citedin{Ch.~\hyperref[ch:signals]{2}}
\item A.~V.~Oppenheim and R.~W.~Schafer, \emph{Discrete-Time Signal Processing}, 3rd ed., Pearson, 2010.\citedin{Ch.~\hyperref[ch:dsp]{6}}
\item A.~Papoulis and S.~U.~Pillai, \emph{Probability, Random Variables and Stochastic Processes}, 4th ed., McGraw-Hill, 2002.\citedin{Ch.~\hyperref[ch:noise]{3}; App.~\hyperref[appA]{A}}
\item B.~W.~Parkinson and J.~J.~Spilker Jr.\ (eds.), \emph{Global Positioning System: Theory and Applications}, 2 vols., AIAA, 1996.\citedin{Ch.~\hyperref[ch:spreadspectrum]{18}}
\item J.~D.~Parsons, \emph{The Mobile Radio Propagation Channel}, 2nd ed., Wiley, 2000.\citedin{Ch.~\hyperref[ch:channels]{11}}
\item C.~R.~Paul, \emph{Introduction to Electromagnetic Compatibility}, 2nd ed., Wiley, 2006.\citedin{Ch.~\hyperref[ch:signals]{2}}
\item A.~Paulraj, R.~Nabar and D.~Gore, \emph{Introduction to Space--Time Wireless Communications}, Cambridge University Press, 2003.\citedin{Ch.~\hyperref[ch:mimo]{19}}
\item E.~Perahia and R.~Stacey, \emph{Next Generation Wireless LANs: 802.11n and 802.11ac}, 2nd ed., Cambridge University Press, 2013.\citedin{Ch.~\hyperref[ch:wifi]{22}}
\item D.~M.~Pozar, \emph{Microwave and RF Design of Wireless Systems}, Wiley, 2001.\citedin{Ch.~\hyperref[ch:sdr]{7}}
\item D.~M.~Pozar, \emph{Microwave Engineering}, 4th ed., Wiley, 2012.\citedin{Ch.~\hyperref[ch:noise]{3}}
\item T.~Pratt and J.~E.~Allnutt, \emph{Satellite Communications}, 3rd ed., Wiley, 2020.\citedin{Ch.~\hyperref[ch:satellite]{23}}
\item J.~G.~Proakis and M.~Salehi, \emph{Digital Communications}, 5th ed., McGraw-Hill, 2008.\citedin{Chs.~\hyperref[ch:baseband]{8}, \hyperref[ch:modulation]{9}, \hyperref[ch:equalization]{12}}
\item L.~R.~Rabiner and R.~W.~Schafer, \emph{Theory and Applications of Digital Speech Processing}, Pearson, 2011.\citedin{Ch.~\hyperref[ch:sourcecoding]{16}}
\item T.~S.~Rappaport, \emph{Wireless Communications: Principles and Practice}, 2nd ed., Prentice Hall, 2002.\citedin{Chs.~\hyperref[ch:channels]{11}, \hyperref[ch:multipleaccess]{20}, \hyperref[ch:cellular]{21}}
\item B.~Razavi, \emph{RF Microelectronics}, 2nd ed., Prentice Hall, 2012.\citedin{Ch.~\hyperref[ch:sdr]{7}}
\item J.~H.~Reed, \emph{Software Radio: A Modern Approach to Radio Engineering}, Prentice Hall, 2002.\citedin{Ch.~\hyperref[ch:sdr]{7}}
\item M.~Rice, \emph{Digital Communications: A Discrete-Time Approach}, Pearson, 2009.\citedin{Ch.~\hyperref[ch:sync]{10}}
\item T.~Richardson and R.~Urbanke, \emph{Modern Coding Theory}, Cambridge University Press, 2008.\citedin{Ch.~\hyperref[ch:moderncodes]{15}}
\item U.~L.~Rohde, J.~C.~Whitaker and H.~Zahnd, \emph{Communications Receivers: Principles and Design}, 4th ed., McGraw-Hill, 2017.\citedin{Ch.~\hyperref[ch:analog]{4}}
\item W.~E.~Ryan and S.~Lin, \emph{Channel Codes: Classical and Modern}, Cambridge University Press, 2009.\citedin{Ch.~\hyperref[ch:moderncodes]{15}}
\item S.~R.~Saunders and A.~Arag\'on-Zavala, \emph{Antennas and Propagation for Wireless Communication Systems}, 2nd ed., Wiley, 2007.\citedin{Ch.~\hyperref[ch:channels]{11}}
\item A.~H.~Sayed, \emph{Adaptive Filters}, Wiley, 2008.\citedin{Ch.~\hyperref[ch:equalization]{12}}
\item K.~Sayood, \emph{Introduction to Data Compression}, 5th ed., Morgan Kaufmann, 2017.\citedin{Ch.~\hyperref[ch:sourcecoding]{16}}
\item R.~Schreier and G.~C.~Temes, \emph{Understanding Delta-Sigma Data Converters}, Wiley-IEEE, 2005; 2nd ed., with S.~Pavan, 2017.\citedin{Ch.~\hyperref[ch:pcm]{5}}
\item S.~Sesia, I.~Toufik and M.~Baker (eds.), \emph{LTE---The UMTS Long Term Evolution: From Theory to Practice}, 2nd ed., Wiley, 2011.\citedin{Ch.~\hyperref[ch:cellular]{21}}
\item M.~K.~Simon, J.~K.~Omura, R.~A.~Scholtz and B.~K.~Levitt, \emph{Spread Spectrum Communications Handbook}, revised ed., McGraw-Hill, 1994.\citedin{Ch.~\hyperref[ch:spreadspectrum]{18}}
\item M.~K.~Simon, S.~M.~Hinedi and W.~C.~Lindsey, \emph{Digital Communication Techniques: Signal Design and Detection}, Prentice Hall, 1995.\citedin{Ch.~\hyperref[ch:modulation]{9}}
\item M.~K.~Simon and M.-S.~Alouini, \emph{Digital Communication over Fading Channels}, 2nd ed., Wiley, 2005.\citedin{Ch.~\hyperref[ch:modulation]{9}; App.~\hyperref[appA]{A}}
\item J.~Soni and R.~Goodman, \emph{A Mind at Play: How Claude Shannon Invented the Information Age}, Simon \& Schuster, 2017.\citedin{Chs.~\hyperref[ch:history]{1}, \hyperref[ch:infotheory]{13}}
\item T.~Standage, \emph{The Victorian Internet}, Walker, 1998.\citedin{Ch.~\hyperref[ch:history]{1}}
\item T.~Starr, J.~M.~Cioffi and P.~J.~Silverman, \emph{Understanding Digital Subscriber Line Technology}, Prentice Hall, 1999.\citedin{Chs.~\hyperref[ch:ofdm]{17}, \hyperref[ch:wireline]{24}}
\item T.~Starr, M.~Sorbara, J.~M.~Cioffi and P.~J.~Silverman, \emph{DSL Advances}, 2003.\citedin{Ch.~\hyperref[ch:wireline]{24}}
\item D.~S.~Taubman and M.~W.~Marcellin, \emph{JPEG2000: Image Compression Fundamentals, Standards and Practice}, Kluwer, 2002.\citedin{Ch.~\hyperref[ch:sourcecoding]{16}}
\item F.~E.~Terman, \emph{Radio Engineers' Handbook}, McGraw-Hill, 1943.\citedin{Ch.~\hyperref[ch:analog]{4}}
\item K.~Townsend, C.~Cuf\'i, Akiba and R.~Davidson, \emph{Getting Started with Bluetooth Low Energy}, O'Reilly, 2014.\citedin{Ch.~\hyperref[ch:wifi]{22}}
\item D.~Tse and P.~Viswanath, \emph{Fundamentals of Wireless Communication}, Cambridge University Press, 2005.\citedin{Chs.~\hyperref[ch:channels]{11}, \hyperref[ch:infotheory]{13}, \hyperref[ch:mimo]{19}, \hyperref[ch:multipleaccess]{20}}
\item P.~P.~Vaidyanathan, \emph{Multirate Systems and Filter Banks}, Prentice Hall, 1993.\citedin{Ch.~\hyperref[ch:dsp]{6}}
\item R.~van Nee and R.~Prasad, \emph{OFDM for Wireless Multimedia Communications}, Artech House, 2000.\citedin{Ch.~\hyperref[ch:ofdm]{17}}
\item H.~L.~Van Trees, \emph{Detection, Estimation, and Modulation Theory, Part~I}, Wiley, 1968 (reissued 2001).\citedin{Ch.~\hyperref[ch:noise]{3}}
\item H.~L.~Van Trees, \emph{Optimum Array Processing} (Part~IV of \emph{Detection, Estimation, and Modulation Theory}), Wiley, 2002.\citedin{Ch.~\hyperref[ch:mimo]{19}}
\item S.~Verd\'u, \emph{Multiuser Detection}, Cambridge University Press, 1998.\citedin{Ch.~\hyperref[ch:spreadspectrum]{18}}
\item A.~J.~Viterbi, \emph{Principles of Coherent Communication}, McGraw-Hill, 1966.\citedin{Ch.~\hyperref[ch:sync]{10}}
\item A.~J.~Viterbi, \emph{CDMA: Principles of Spread Spectrum Communication}, Addison-Wesley, 1995.\citedin{Chs.~\hyperref[ch:spreadspectrum]{18}, \hyperref[ch:cellular]{21}}
\item S.~B.~Wicker and V.~K.~Bhargava (eds.), \emph{Reed--Solomon Codes and Their Applications}, IEEE Press, 1994.\citedin{Ch.~\hyperref[ch:classiccodes]{14}}
\item B.~Widrow and S.~D.~Stearns, \emph{Adaptive Signal Processing}, Prentice Hall, 1985.\citedin{Ch.~\hyperref[ch:equalization]{12}}
\item B.~Widrow and I.~Koll\'ar, \emph{Quantization Noise: Roundoff Error in Digital Computation, Signal Processing, Control, and Communications}, Cambridge University Press, 2008.\citedin{Ch.~\hyperref[ch:pcm]{5}}
\item J.~M.~Wozencraft and I.~M.~Jacobs, \emph{Principles of Communication Engineering}, Wiley, 1965.\citedin{Ch.~\hyperref[ch:modulation]{9}}
\item J.~H.~Yuen (ed.), \emph{Deep Space Telecommunications Systems Engineering}, JPL, 1982 (free online), and the JPL DESCANSO monograph series.\citedin{Chs.~\hyperref[ch:noise]{3}, \hyperref[ch:satellite]{23}}
\end{itemize}
\endgroup

\section*{Classic and research papers}
\addcontentsline{toc}{section}{Classic and research papers}
\markright{Classic and research papers}
\begingroup\small
\begin{itemize}[leftmargin=1.5em,itemindent=-1.5em,itemsep=1.5pt plus 1pt,parsep=0pt,topsep=4pt,label={}]\rightskip=0pt plus 3em\relax
\item A.~A.~Abidi, ``Direct-conversion radio transceivers for digital communications,'' \emph{IEEE Journal of Solid-State Circuits}, vol.~30, no.~12, 1995.\citedin{Ch.~\hyperref[ch:sdr]{7}}
\item N.~Abramson, ``The ALOHA system---another alternative for computer communications,'' \emph{Proceedings of the AFIPS Fall Joint Computer Conference}, 1970.\citedin{Ch.~\hyperref[ch:multipleaccess]{20}}
\item S.~M.~Alamouti, ``A simple transmit diversity technique for wireless communications,'' \emph{IEEE Journal on Selected Areas in Communications}, vol.~16, no.~8, 1998.\citedin{Ch.~\hyperref[ch:mimo]{19}}
\item J.~B.~Anderson, F.~Rusek and V.~\"Owall, ``Faster-than-Nyquist signaling,'' \emph{Proceedings of the IEEE}, vol.~101, no.~8, 2013.\citedin{Ch.~\hyperref[ch:baseband]{8}}
\item R.~Andraka, ``A survey of CORDIC algorithms for FPGA based computers,'' \emph{Proceedings of the ACM/SIGDA FPGA~'98}, 1998.\citedin{Ch.~\hyperref[ch:dsp]{6}}
\item J.~G.~Andrews, F.~Baccelli and R.~K.~Ganti, ``A tractable approach to coverage and rate in cellular networks,'' \emph{IEEE Transactions on Communications}, vol.~59, no.~11, 2011.\citedin{Ch.~\hyperref[ch:multipleaccess]{20}}
\item E.~Ar\i kan, ``Channel polarization: a method for constructing capacity-achieving codes for symmetric binary-input memoryless channels,'' \emph{IEEE Transactions on Information Theory}, vol.~55, no.~7, 2009.\citedin{Ch.~\hyperref[ch:moderncodes]{15}}
\item E.~H.~Armstrong, ``A method of reducing disturbances in radio signaling by a system of frequency modulation,'' \emph{Proceedings of the IRE}, vol.~24, no.~5, pp.~689--740, 1936.\citedin{Ch.~\hyperref[ch:analog]{4}}
\item L.~R.~Bahl, J.~Cocke, F.~Jelinek and J.~Raviv, ``Optimal decoding of linear codes for minimizing symbol error rate,'' \emph{IEEE Transactions on Information Theory}, vol.~20, no.~2, 1974.\citedin{Ch.~\hyperref[ch:moderncodes]{15}}
\item J.~Ball\'e, V.~Laparra and E.~P.~Simoncelli, ``End-to-end optimized image compression,'' \emph{International Conference on Learning Representations} (ICLR), 2017.\citedin{Ch.~\hyperref[ch:sourcecoding]{16}}
\item P.~Banelli \emph{et al.}, ``Modulation formats and waveforms for 5G networks: who will be the heir of OFDM?,'' \emph{IEEE Signal Processing Magazine}, vol.~31, no.~6, 2014.\citedin{Ch.~\hyperref[ch:ofdm]{17}}
\item E.~Barkan, E.~Biham and N.~Keller, ``Instant ciphertext-only cryptanalysis of GSM encrypted communication,'' \emph{Journal of Cryptology}, vol.~21, no.~3, 2008 (first presented at CRYPTO~2003).\citedin{Ch.~\hyperref[ch:cellular]{21}}
\item M.~G.~Bellanger, G.~Bonnerot and M.~Coudreuse, ``Digital filtering by polyphase network: application to sample-rate alteration and filter banks,'' \emph{IEEE Transactions on Acoustics, Speech, and Signal Processing}, vol.~24, no.~2, 1976.\citedin{Ch.~\hyperref[ch:dsp]{6}}
\item P.~A.~Bello, ``Characterization of randomly time-variant linear channels,'' \emph{IEEE Transactions on Communications Systems}, vol.~CS-11, 1963.\citedin{Ch.~\hyperref[ch:channels]{11}}
\item A.~Bemani, N.~Ksairi and M.~Kountouris, ``Affine frequency division multiplexing for next generation wireless communications,'' \emph{IEEE Transactions on Wireless Communications}, vol.~22, no.~11, 2023.\citedin{Ch.~\hyperref[ch:sixg]{25}}
\item C.~H.~Bennett and G.~Brassard, ``Quantum cryptography: public key distribution and coin tossing,'' \emph{Proceedings of the IEEE International Conference on Computers, Systems and Signal Processing}, Bangalore, 1984.\citedin{Ch.~\hyperref[ch:sixg]{25}}
\item C.~Berrou, A.~Glavieux and P.~Thitimajshima, ``Near Shannon limit error-correcting coding and decoding: Turbo-codes,'' \emph{Proceedings of IEEE ICC}, Geneva, pp.~1064--1070, 1993.\citedin{Ch.~\hyperref[ch:moderncodes]{15}}
\item G.~Bianchi, ``Performance analysis of the IEEE 802.11 distributed coordination function,'' \emph{IEEE Journal on Selected Areas in Communications}, vol.~18, no.~3, 2000.\citedin{Ch.~\hyperref[ch:multipleaccess]{20}}
\item J.~A.~C.~Bingham, ``Multicarrier modulation for data transmission: an idea whose time has come,'' \emph{IEEE Communications Magazine}, vol.~28, no.~5, 1990.\citedin{Ch.~\hyperref[ch:ofdm]{17}}
\item V.~Bioglio, C.~Condo and I.~Land, ``Design of polar codes in 5G New Radio,'' \emph{IEEE Communications Surveys \& Tutorials}, vol.~23, no.~1, 2021.\citedin{Ch.~\hyperref[ch:moderncodes]{15}}
\item E.~Bj\"ornson, L.~Sanguinetti, H.~Wymeersch, J.~Hoydis and T.~L.~Marzetta, ``Massive MIMO is a reality---What is next? Five promising research directions for antenna arrays,'' \emph{Digital Signal Processing}, vol.~94, 2019.\citedin{Ch.~\hyperref[ch:sixg]{25}}
\item E.~Bj\"ornson, \"O.~\"Ozdogan and E.~G.~Larsson, ``Intelligent reflecting surface versus decode-and-forward: how large surfaces are needed to beat relaying?,'' \emph{IEEE Wireless Communications Letters}, vol.~9, no.~2, 2020.\citedin{Ch.~\hyperref[ch:sixg]{25}}
\item H.~S.~Black, ``Inventing the negative feedback amplifier,'' \emph{IEEE Spectrum}, December 1977.\citedin{Ch.~\hyperref[ch:history]{1}}
\item G.~B\"ocherer, F.~Steiner and P.~Schulte, ``Bandwidth efficient and rate-matched low-density parity-check coded modulation,'' \emph{IEEE Transactions on Communications}, vol.~63, no.~12, 2015.\citedin{Ch.~\hyperref[ch:modulation]{9}}
\item B.~Bross \emph{et al.}, ``Overview of the Versatile Video Coding (VVC) standard and its applications,'' \emph{IEEE Transactions on Circuits and Systems for Video Technology}, vol.~31, no.~10, 2021.\citedin{Ch.~\hyperref[ch:sourcecoding]{16}}
\item G.~Caire, G.~Taricco and E.~Biglieri, ``Bit-interleaved coded modulation,'' \emph{IEEE Transactions on Information Theory}, vol.~44, no.~3, pp.~927--946, 1998.\citedin{Chs.~\hyperref[ch:modulation]{9}, \hyperref[ch:infotheory]{13}, \hyperref[ch:ofdm]{17}}
\item J.~R.~Carson, ``Notes on the theory of modulation,'' \emph{Proceedings of the IRE}, vol.~10, no.~1, pp.~57--64, 1922.\citedin{Ch.~\hyperref[ch:analog]{4}}
\item R.~W.~Chang, ``Synthesis of band-limited orthogonal signals for multichannel data transmission,'' \emph{Bell System Technical Journal}, vol.~45, 1966.\citedin{Ch.~\hyperref[ch:ofdm]{17}}
\item M.~Chiani, D.~Dardari and M.~K.~Simon, ``New exponential bounds and approximations for the computation of error probability in fading channels,'' \emph{IEEE Transactions on Wireless Communications}, vol.~2, no.~4, 2003.\citedin{App.~\hyperref[appA]{A}}
\item R.~D.~Cideciyan, F.~Dolivo, R.~Hermann, W.~Hirt and W.~Schott, ``A PRML system for digital magnetic recording,'' \emph{IEEE Journal on Selected Areas in Communications}, vol.~10, no.~1, 1992.\citedin{Ch.~\hyperref[ch:baseband]{8}}
\item J.~M.~Cioffi, G.~P.~Dudevoir, M.~V.~Eyuboglu and G.~D.~Forney, Jr., ``MMSE decision-feedback equalizers and coding,'' Parts~I and~II, \emph{IEEE Transactions on Communications}, vol.~43, no.~10, 1995.\citedin{Ch.~\hyperref[ch:equalization]{12}}
\item A.~C.~Clarke, ``Extra-terrestrial relays,'' \emph{Wireless World}, October 1945.\citedin{Chs.~\hyperref[ch:history]{1}, \hyperref[ch:satellite]{23}}
\item R.~H.~Clarke, ``A statistical theory of mobile-radio reception,'' \emph{Bell System Technical Journal}, vol.~47, 1968.\citedin{Ch.~\hyperref[ch:channels]{11}}
\item J.~W.~Cooley and J.~W.~Tukey, ``An algorithm for the machine calculation of complex Fourier series,'' \emph{Mathematics of Computation}, vol.~19, no.~90, pp.~297--301, 1965.\citedin{Ch.~\hyperref[ch:signals]{2}}
\item J.~P.~Costas, ``Synchronous communications,'' \emph{Proceedings of the IRE}, vol.~44, no.~12, 1956.\citedin{Ch.~\hyperref[ch:analog]{4}}
\item D.~J.~Costello, Jr.\ and G.~D.~Forney, Jr., ``Channel coding: the road to channel capacity,'' \emph{Proceedings of the IEEE}, vol.~95, no.~6, pp.~1150--1177, 2007.\citedin{Chs.~\hyperref[ch:infotheory]{13}, \hyperref[ch:classiccodes]{14}}
\item J.~W.~Craig, ``A new, simple and exact result for calculating the probability of error for two-dimensional signal constellations,'' \emph{Proceedings of IEEE MILCOM}, 1991.\citedin{App.~\hyperref[appA]{A}}
\item R.~De~Gaudenzi, A.~Guill\'en~i~F\`abregas and A.~Martinez, ``Turbo-coded APSK modulations design for satellite broadband communications,'' \emph{International Journal of Satellite Communications and Networking}, 2006.\citedin{Ch.~\hyperref[ch:satellite]{23}}
\item C.~Deng \emph{et al.}, ``IEEE 802.11be Wi-Fi 7: New challenges and opportunities,'' \emph{IEEE Communications Surveys \& Tutorials}, vol.~22, no.~4, 2020.\citedin{Ch.~\hyperref[ch:wifi]{22}}
\item W.~H.~Doherty, ``A new high efficiency power amplifier for modulated waves,'' \emph{Proceedings of the IRE}, vol.~24, no.~9, 1936.\citedin{Ch.~\hyperref[ch:sdr]{7}}
\item O.~Edfors, M.~Sandell, J.-J.~van de Beek, S.~K.~Wilson and P.~O.~B\"orjesson, ``OFDM channel estimation by singular value decomposition,'' \emph{IEEE Transactions on Communications}, vol.~46, no.~7, 1998.\citedin{Ch.~\hyperref[ch:ofdm]{17}}
\item L.~Erup, F.~M.~Gardner and R.~A.~Harris, ``Interpolation in digital modems---Part~II: Implementation and performance,'' \emph{IEEE Transactions on Communications}, vol.~41, no.~6, 1993.\citedin{Ch.~\hyperref[ch:sync]{10}}
\item R.-J.~Essiambre, G.~Kramer, P.~J.~Winzer, G.~J.~Foschini and B.~Goebel, ``Capacity limits of optical fiber networks,'' \emph{Journal of Lightwave Technology}, vol.~28, no.~4, 2010.\citedin{Ch.~\hyperref[ch:wireline]{24}}
\item D.~Falconer, S.~L.~Ariyavisitakul, A.~Benyamin-Seeyar and B.~Eidson, ``Frequency domain equalization for single-carrier broadband wireless systems,'' \emph{IEEE Communications Magazine}, vol.~40, no.~4, 2002.\citedin{Ch.~\hyperref[ch:equalization]{12}}
\item B.~Farhang-Boroujeny, ``OFDM versus filter bank multicarrier,'' \emph{IEEE Signal Processing Magazine}, vol.~28, no.~3, 2011.\citedin{Ch.~\hyperref[ch:ofdm]{17}}
\item G.~D.~Forney, Jr., ``Maximum-likelihood sequence estimation of digital sequences in the presence of intersymbol interference,'' \emph{IEEE Transactions on Information Theory}, vol.~18, no.~3, 1972.\citedin{Ch.~\hyperref[ch:equalization]{12}}
\item G.~D.~Forney, Jr., ``The Viterbi algorithm,'' \emph{Proceedings of the IEEE}, vol.~61, no.~3, 1973.\citedin{Ch.~\hyperref[ch:classiccodes]{14}}
\item G.~D.~Forney, Jr.\ and G.~Ungerboeck, ``Modulation and coding for linear Gaussian channels,'' \emph{IEEE Transactions on Information Theory}, vol.~44, no.~6, pp.~2384--2415, 1998.\citedin{Chs.~\hyperref[ch:modulation]{9}, \hyperref[ch:infotheory]{13}}
\item G.~J.~Foschini and M.~J.~Gans, ``On limits of wireless communications in a fading environment when using multiple antennas,'' \emph{Wireless Personal Communications}, vol.~6, 1998.\citedin{Ch.~\hyperref[ch:mimo]{19}}
\item H.~T.~Friis, ``Noise figures of radio receivers,'' \emph{Proceedings of the IRE}, vol.~32, pp.~419--422, 1944.\citedin{Ch.~\hyperref[ch:noise]{3}}
\item H.~T.~Friis, ``A note on a simple transmission formula,'' \emph{Proceedings of the IRE}, vol.~34, pp.~254--256, 1946.\citedin{Ch.~\hyperref[ch:noise]{3}}
\item D.~Gabor, ``Theory of communication,'' \emph{Journal of the IEE, Part~III}, vol.~93, no.~26, pp.~429--457, 1946.\citedin{Ch.~\hyperref[ch:signals]{2}}
\item F.~M.~Gardner, ``A BPSK/QPSK timing-error detector for sampled receivers,'' \emph{IEEE Transactions on Communications}, vol.~34, no.~5, 1986.\citedin{Ch.~\hyperref[ch:sync]{10}}
\item W.~A.~Gardner, ``Exploitation of spectral redundancy in cyclostationary signals,'' \emph{IEEE Signal Processing Magazine}, vol.~8, no.~2, pp.~14--36, 1991.\citedin{Ch.~\hyperref[ch:noise]{3}}
\item F.~M.~Gardner, ``Interpolation in digital modems---Part~I: Fundamentals,'' \emph{IEEE Transactions on Communications}, vol.~41, no.~3, 1993.\citedin{Ch.~\hyperref[ch:sync]{10}}
\item M.~Gastpar, B.~Rimoldi and M.~Vetterli, ``To code, or not to code: lossy source--channel communication revisited,'' \emph{IEEE Transactions on Information Theory}, vol.~49, no.~5, 2003.\citedin{Ch.~\hyperref[ch:sourcecoding]{16}}
\item K.~S.~Gilhousen, I.~M.~Jacobs, R.~Padovani, A.~J.~Viterbi, L.~A.~Weaver and C.~E.~Wheatley, ``On the capacity of a cellular CDMA system,'' \emph{IEEE Transactions on Vehicular Technology}, vol.~40, no.~2, 1991.\citedin{Ch.~\hyperref[ch:spreadspectrum]{18}}
\item G.~Ginis and J.~M.~Cioffi, ``Vectored transmission for digital subscriber line systems,'' \emph{IEEE Journal on Selected Areas in Communications}, vol.~20, no.~5, 2002.\citedin{Ch.~\hyperref[ch:wireline]{24}}
\item R.~M.~Gray, ``Toeplitz and circulant matrices: a review,'' \emph{Foundations and Trends in Communications and Information Theory}, vol.~2, no.~3, 2006.\citedin{App.~\hyperref[appA]{A}}
\item D.~G\"und\"uz \emph{et al.}, ``Beyond transmitting bits: context, semantics, and task-oriented communications,'' \emph{IEEE Journal on Selected Areas in Communications}, vol.~41, no.~1, 2023.\citedin{Ch.~\hyperref[ch:sixg]{25}}
\item R.~Hadani \emph{et al.}, ``Orthogonal time frequency space modulation,'' \emph{Proceedings of IEEE WCNC}, 2017.\citedin{Ch.~\hyperref[ch:sixg]{25}}
\item R.~W.~Hamming, ``Error detecting and error correcting codes,'' \emph{Bell System Technical Journal}, vol.~29, no.~2, 1950.\citedin{Ch.~\hyperref[ch:classiccodes]{14}}
\item S.~H.~Han and J.~H.~Lee, ``An overview of peak-to-average power ratio reduction techniques for multicarrier transmission,'' \emph{IEEE Wireless Communications}, vol.~12, no.~2, 2005.\citedin{Ch.~\hyperref[ch:ofdm]{17}}
\item F.~J.~Harris, ``On the use of windows for harmonic analysis with the discrete Fourier transform,'' \emph{Proceedings of the IEEE}, vol.~66, no.~1, pp.~51--83, 1978.\citedin{Ch.~\hyperref[ch:signals]{2}; App.~\hyperref[appA]{A}}
\item f.~j.~harris and M.~Rice, ``Multirate digital filters for symbol timing synchronization in software defined radios,'' \emph{IEEE Journal on Selected Areas in Communications}, vol.~19, no.~12, 2001.\citedin{Ch.~\hyperref[ch:sync]{10}}
\item R.~V.~L.~Hartley, ``Transmission of information,'' \emph{Bell System Technical Journal}, vol.~7, 1928.\citedin{Ch.~\hyperref[ch:history]{1}}
\item M.~Hata, ``Empirical formula for propagation loss in land mobile radio services,'' \emph{IEEE Transactions on Vehicular Technology}, vol.~VT-29, 1980.\citedin{Ch.~\hyperref[ch:channels]{11}}
\item R.~W.~Heath Jr., N.~Gonz\'alez-Prelcic, S.~Rangan, W.~Roh and A.~M.~Sayeed, ``An overview of signal processing techniques for millimeter wave MIMO systems,'' \emph{IEEE Journal of Selected Topics in Signal Processing}, vol.~10, no.~3, 2016.\citedin{Ch.~\hyperref[ch:mimo]{19}}
\item M.~T.~Heideman, D.~H.~Johnson and C.~S.~Burrus, ``Gauss and the history of the fast Fourier transform,'' \emph{IEEE ASSP Magazine}, vol.~1, no.~4, pp.~14--21, 1984.\citedin{Ch.~\hyperref[ch:signals]{2}}
\item M.~Heusse, F.~Rousseau, G.~Berger-Sabbatel and A.~Duda, ``Performance anomaly of 802.11b,'' \emph{Proceedings of IEEE INFOCOM}, 2003.\citedin{Ch.~\hyperref[ch:wifi]{22}}
\item E.~B.~Hogenauer, ``An economical class of digital filters for decimation and interpolation,'' \emph{IEEE Transactions on Acoustics, Speech, and Signal Processing}, vol.~29, no.~2, pp.~155--162, 1981.\citedin{Ch.~\hyperref[ch:dsp]{6}}
\item H.~Inose, Y.~Yasuda and J.~Murakami, ``A telemetering system by code modulation: $\Delta$--$\Sigma$ modulation,'' \emph{IRE Transactions on Space Electronics and Telemetry}, vol.~SET-8, 1962.\citedin{Ch.~\hyperref[ch:pcm]{5}}
\item J.~B.~Johnson, ``Thermal agitation of electricity in conductors,'' \emph{Physical Review}, vol.~32, pp.~97--109, 1928.\citedin{Ch.~\hyperref[ch:noise]{3}}
\item C.~R.~Johnson, Jr.\ \emph{et al.}, ``Blind equalization using the constant modulus criterion: a review,'' \emph{Proceedings of the IEEE}, vol.~86, no.~10, 1998.\citedin{Ch.~\hyperref[ch:equalization]{12}}
\item P.~Kabal and S.~Pasupathy, ``Partial-response signaling,'' \emph{IEEE Transactions on Communications}, vol.~23, no.~9, 1975.\citedin{Ch.~\hyperref[ch:baseband]{8}}
\item K.~C.~Kao and G.~A.~Hockham, ``Dielectric-fibre surface waveguides for optical frequencies,'' \emph{Proceedings of the IEE}, vol.~113, no.~7, 1966.\citedin{Ch.~\hyperref[ch:wireline]{24}}
\item S.~Kay, ``A fast and accurate single frequency estimator,'' \emph{IEEE Transactions on Acoustics, Speech, and Signal Processing}, vol.~37, no.~12, 1989.\citedin{Ch.~\hyperref[ch:sync]{10}}
\item F.~P.~Kelly, A.~K.~Maulloo and D.~K.~H.~Tan, ``Rate control for communication networks: shadow prices, proportional fairness and stability,'' \emph{Journal of the Operational Research Society}, vol.~49, 1998.\citedin{Ch.~\hyperref[ch:multipleaccess]{20}}
\item E.~Khorov, A.~Kiryanov, A.~Lyakhov and G.~Bianchi, ``A tutorial on IEEE 802.11ax high efficiency WLANs,'' \emph{IEEE Communications Surveys \& Tutorials}, vol.~21, no.~1, 2019.\citedin{Ch.~\hyperref[ch:wifi]{22}}
\item K.~Kikuchi, ``Fundamentals of coherent optical fiber communications,'' \emph{Journal of Lightwave Technology}, vol.~34, no.~1, 2016.\citedin{Ch.~\hyperref[ch:wireline]{24}}
\item L.~Kleinrock and F.~A.~Tobagi, ``Packet switching in radio channels: Part~I---Carrier sense multiple-access modes and their throughput-delay characteristics,'' \emph{IEEE Transactions on Communications}, vol.~23, no.~12, 1975.\citedin{Ch.~\hyperref[ch:multipleaccess]{20}}
\item R.~Koetter, A.~C.~Singer and M.~T\"uchler, ``Turbo equalization,'' \emph{IEEE Signal Processing Magazine}, vol.~21, no.~1, 2004.\citedin{Ch.~\hyperref[ch:equalization]{12}}
\item P.~Koopman, ``32-bit cyclic redundancy codes for Internet applications,'' \emph{Proceedings of the International Conference on Dependable Systems and Networks}, 2002.\citedin{Ch.~\hyperref[ch:classiccodes]{14}}
\item F.~R.~Kschischang and S.~Pasupathy, ``Optimal nonuniform signaling for Gaussian channels,'' \emph{IEEE Transactions on Information Theory}, vol.~39, no.~3, 1993.\citedin{Ch.~\hyperref[ch:modulation]{9}}
\item S.~Kudekar, T.~Richardson and R.~Urbanke, ``Threshold saturation via spatial coupling: why convolutional LDPC ensembles perform so well over the BEC,'' \emph{IEEE Transactions on Information Theory}, vol.~57, no.~2, 2011.\citedin{Ch.~\hyperref[ch:moderncodes]{15}}
\item D.~B.~Leeson, ``A simple model of feedback oscillator noise spectrum,'' \emph{Proceedings of the IEEE}, vol.~54, no.~2, 1966.\citedin{Ch.~\hyperref[ch:sdr]{7}}
\item Y.~Li, L.~J.~Cimini and N.~R.~Sollenberger, ``Robust channel estimation for OFDM systems with rapid dispersive fading channels,'' \emph{IEEE Transactions on Communications}, vol.~46, no.~7, 1998.\citedin{Ch.~\hyperref[ch:ofdm]{17}}
\item Y.~Li and L.~J.~Cimini, ``Bounds on the interchannel interference of OFDM in time-varying impairments,'' \emph{IEEE Transactions on Communications}, vol.~49, no.~3, 2001.\citedin{Ch.~\hyperref[ch:ofdm]{17}}
\item S.~P.~Lipshitz, R.~A.~Wannamaker and J.~Vanderkooy, ``Quantization and dither: a theoretical survey,'' \emph{Journal of the Audio Engineering Society}, vol.~40, no.~5, 1992.\citedin{Ch.~\hyperref[ch:pcm]{5}}
\item V.~Liu, A.~Parks, V.~Talla, S.~Gollakota, D.~Wetherall and J.~R.~Smith, ``Ambient backscatter: Wireless communication out of thin air,'' \emph{Proceedings of ACM SIGCOMM}, 2013.\citedin{Ch.~\hyperref[ch:wifi]{22}}
\item F.~Liu \emph{et al.}, ``Integrated sensing and communications: toward dual-functional wireless networks for 6G and beyond,'' \emph{IEEE Journal on Selected Areas in Communications}, vol.~40, no.~6, 2022.\citedin{Ch.~\hyperref[ch:sixg]{25}}
\item R.~W.~Lucky, ``Automatic equalization for digital communication,'' \emph{Bell System Technical Journal}, vol.~44, no.~4, 1965.\citedin{Ch.~\hyperref[ch:equalization]{12}}
\item M.~Luise and R.~Reggiannini, ``Carrier frequency recovery in all-digital modems for burst-mode transmissions,'' \emph{IEEE Transactions on Communications}, vol.~43, no.~2/3/4, 1995.\citedin{Ch.~\hyperref[ch:sync]{10}}
\item V.~H.~MacDonald, ``Advanced Mobile Phone Service: The cellular concept,'' \emph{Bell System Technical Journal}, vol.~58, no.~1, pp.~15--41, 1979.\citedin{Chs.~\hyperref[ch:multipleaccess]{20}, \hyperref[ch:cellular]{21}}
\item D.~J.~C.~MacKay, ``Good error-correcting codes based on very sparse matrices,'' \emph{IEEE Transactions on Information Theory}, vol.~45, no.~2, 1999.\citedin{Ch.~\hyperref[ch:moderncodes]{15}}
\item T.~L.~Marzetta, ``Noncooperative cellular wireless with unlimited numbers of base station antennas,'' \emph{IEEE Transactions on Wireless Communications}, vol.~9, no.~11, 2010.\citedin{Ch.~\hyperref[ch:mimo]{19}}
\item J.~L.~Massey, ``Optimum frame synchronization,'' \emph{IEEE Transactions on Communications}, vol.~20, no.~2, 1972.\citedin{Ch.~\hyperref[ch:sync]{10}}
\item J.~E.~Mazo, ``Faster-than-Nyquist signaling,'' \emph{Bell System Technical Journal}, vol.~54, no.~8, 1975.\citedin{Ch.~\hyperref[ch:baseband]{8}}
\item J.~H.~McClellan, T.~W.~Parks and L.~R.~Rabiner, ``A computer program for designing optimum FIR linear phase digital filters,'' \emph{IEEE Transactions on Audio and Electroacoustics}, vol.~21, no.~6, 1973.\citedin{Ch.~\hyperref[ch:dsp]{6}}
\item R.~M.~Metcalfe and D.~R.~Boggs, ``Ethernet: Distributed packet switching for local computer networks,'' \emph{Communications of the ACM}, vol.~19, no.~7, 1976.\citedin{Ch.~\hyperref[ch:multipleaccess]{20}}
\item J.~Mitola, ``The software radio architecture,'' \emph{IEEE Communications Magazine}, vol.~33, no.~5, 1995.\citedin{Ch.~\hyperref[ch:sdr]{7}}
\item J.~Mo and J.~Walrand, ``Fair end-to-end window-based congestion control,'' \emph{IEEE/ACM Transactions on Networking}, vol.~8, no.~5, 2000.\citedin{Ch.~\hyperref[ch:multipleaccess]{20}}
\item A.~Morello and V.~Mignone, ``DVB-S2: The second generation standard for satellite broad-band services,'' \emph{Proceedings of the IEEE}, vol.~94, no.~1, 2006.\citedin{Ch.~\hyperref[ch:satellite]{23}}
\item D.~R.~Morgan, Z.~Ma, J.~Kim, M.~G.~Zierdt and J.~Pastalan,``A generalized memory polynomial model for digital predistortion of RF power amplifiers,'' \emph{IEEE Transactions on Signal Processing}, vol.~54, no.~10, 2006.\citedin{Ch.~\hyperref[ch:sdr]{7}}
\item K.~H.~Mueller and M.~M\"uller, ``Timing recovery in digital synchronous data receivers,'' \emph{IEEE Transactions on Communications}, vol.~24, no.~5, 1976.\citedin{Ch.~\hyperref[ch:sync]{10}}
\item K.~Murota and K.~Hirade, ``GMSK modulation for digital mobile radio telephony,'' \emph{IEEE Transactions on Communications}, vol.~29, no.~7, 1981.\citedin{Ch.~\hyperref[ch:modulation]{9}}
\item H.~Nyquist, ``Certain factors affecting telegraph speed,'' \emph{Bell System Technical Journal}, vol.~3, 1924.\citedin{Ch.~\hyperref[ch:history]{1}}
\item H.~Nyquist, ``Certain topics in telegraph transmission theory,'' \emph{Transactions of the AIEE}, vol.~47, pp.~617--644, 1928 (reprinted in \emph{Proceedings of the IEEE}, vol.~90, no.~2, 2002).\citedin{Ch.~\hyperref[ch:baseband]{8}}
\item H.~Nyquist, ``Thermal agitation of electric charge in conductors,'' \emph{Physical Review}, vol.~32, pp.~110--113, 1928.\citedin{Ch.~\hyperref[ch:noise]{3}}
\item M.~Oerder and H.~Meyr, ``Digital filter and square timing recovery,'' \emph{IEEE Transactions on Communications}, vol.~36, no.~5, 1988.\citedin{Ch.~\hyperref[ch:sync]{10}}
\item Y.~Okumura \emph{et al.}, ``Field strength and its variability in VHF and UHF land-mobile radio service,'' \emph{Review of the Electrical Communication Laboratory}, vol.~16, 1968.\citedin{Ch.~\hyperref[ch:channels]{11}}
\item B.~M.~Oliver, J.~R.~Pierce and C.~E.~Shannon, ``The philosophy of PCM,'' \emph{Proceedings of the IRE}, vol.~36, no.~11, 1948.\citedin{Ch.~\hyperref[ch:pcm]{5}}
\item T.~O'Shea and J.~Hoydis, ``An introduction to deep learning for the physical layer,'' \emph{IEEE Transactions on Cognitive Communications and Networking}, vol.~3, no.~4, 2017.\citedin{Ch.~\hyperref[ch:sixg]{25}}
\item T.~Painter and A.~Spanias, ``Perceptual coding of digital audio,'' \emph{Proceedings of the IEEE}, vol.~88, no.~4, 2000.\citedin{Ch.~\hyperref[ch:sourcecoding]{16}}
\item T.~W.~Parks and J.~H.~McClellan, ``Chebyshev approximation for nonrecursive digital filters with linear phase,'' \emph{IEEE Transactions on Circuit Theory}, vol.~19, no.~2, 1972.\citedin{Ch.~\hyperref[ch:dsp]{6}}
\item S.~Pasupathy, ``Minimum shift keying: a spectrally efficient modulation,'' \emph{IEEE Communications Magazine}, vol.~17, no.~4, 1979.\citedin{Ch.~\hyperref[ch:modulation]{9}}
\item P.~Poggiolini, ``The GN model of non-linear propagation in uncompensated coherent optical systems,'' \emph{Journal of Lightwave Technology}, vol.~30, no.~24, 2012.\citedin{Ch.~\hyperref[ch:wireline]{24}}
\item T.~Pollet, M.~Van Bladel and M.~Moeneclaey, ``BER sensitivity of OFDM systems to carrier frequency offset and Wiener phase noise,'' \emph{IEEE Transactions on Communications}, vol.~43, 1995.\citedin{Ch.~\hyperref[ch:ofdm]{17}}
\item Y.~Polyanskiy, H.~V.~Poor and S.~Verd\'u, ``Channel coding rate in the finite blocklength regime,'' \emph{IEEE Transactions on Information Theory}, vol.~56, no.~5, pp.~2307--2359, 2010.\citedin{Chs.~\hyperref[ch:infotheory]{13}, \hyperref[ch:moderncodes]{15}}
\item S.~U.~H.~Qureshi, ``Adaptive equalization,'' \emph{Proceedings of the IEEE}, vol.~73, no.~9, 1985.\citedin{Ch.~\hyperref[ch:equalization]{12}}
\item T.~S.~Rappaport \emph{et al.}, ``Millimeter wave mobile communications for 5G cellular: it will work!,'' \emph{IEEE Access}, vol.~1, 2013.\citedin{Ch.~\hyperref[ch:channels]{11}}
\item T.~S.~Rappaport \emph{et al.}, ``Wireless communications and applications above 100\,GHz: opportunities and challenges for 6G and beyond,'' \emph{IEEE Access}, vol.~7, 2019.\citedin{Ch.~\hyperref[ch:sixg]{25}}
\item P.~Raviteja, K.~T.~Phan and Y.~Hong, ``Embedded pilot-aided channel estimation for OTFS in delay--Doppler channels,'' \emph{IEEE Transactions on Vehicular Technology}, vol.~68, no.~5, 2019.\citedin{Ch.~\hyperref[ch:sixg]{25}}
\item I.~S.~Reed and G.~Solomon, ``Polynomial codes over certain finite fields,'' \emph{Journal of the SIAM}, vol.~8, no.~2, 1960.\citedin{Ch.~\hyperref[ch:classiccodes]{14}}
\item S.~O.~Rice, ``Mathematical analysis of random noise,'' \emph{Bell System Technical Journal}, vol.~23, pp.~282--332, 1944, and vol.~24, pp.~46--156, 1945.\citedin{Ch.~\hyperref[ch:noise]{3}}
\item S.~O.~Rice, ``Noise in FM receivers,'' in M.~Rosenblatt (ed.), \emph{Time Series Analysis}, Wiley, 1963.\citedin{Ch.~\hyperref[ch:analog]{4}}
\item T.~Richardson, A.~Shokrollahi and R.~Urbanke, ``Design of capacity-approaching irregular low-density parity-check codes,'' \emph{IEEE Transactions on Information Theory}, vol.~47, no.~2, 2001.\citedin{Ch.~\hyperref[ch:moderncodes]{15}}
\item T.~Richardson and S.~Kudekar, ``Design of low-density parity check codes for 5G New Radio,'' \emph{IEEE Communications Magazine}, vol.~56, no.~3, 2018.\citedin{Ch.~\hyperref[ch:moderncodes]{15}}
\item A.~A.~M.~Saleh, ``Frequency-independent and frequency-dependent nonlinear models of TWT amplifiers,'' \emph{IEEE Transactions on Communications}, vol.~29, no.~11, 1981.\citedin{Ch.~\hyperref[ch:satellite]{23}}
\item D.~V.~Sarwate and M.~B.~Pursley, ``Crosscorrelation properties of pseudorandom and related sequences,'' \emph{Proceedings of the IEEE}, vol.~68, no.~5, 1980.\citedin{Ch.~\hyperref[ch:spreadspectrum]{18}}
\item S.~J.~Savory, ``Digital coherent optical receivers: algorithms and subsystems,'' \emph{IEEE Journal of Selected Topics in Quantum Electronics}, vol.~16, no.~5, 2010.\citedin{Chs.~\hyperref[ch:equalization]{12}, \hyperref[ch:wireline]{24}}
\item T.~M.~Schmidl and D.~C.~Cox, ``Robust frequency and timing synchronization for OFDM,'' \emph{IEEE Transactions on Communications}, vol.~45, no.~12, 1997.\citedin{Ch.~\hyperref[ch:sync]{10}}
\item R.~A.~Scholtz, ``The origins of spread-spectrum communications,'' \emph{IEEE Transactions on Communications}, vol.~30, no.~5, 1982.\citedin{Ch.~\hyperref[ch:spreadspectrum]{18}}
\item C.~E.~Shannon, ``A mathematical theory of communication,'' \emph{Bell System Technical Journal}, vol.~27, pp.~379--423 and 623--656, 1948.\citedin{Chs.~\hyperref[ch:history]{1}, \hyperref[ch:infotheory]{13}}
\item C.~E.~Shannon, ``Communication in the presence of noise,'' \emph{Proceedings of the IRE}, vol.~37, no.~1, pp.~10--21, 1949.\citedin{Ch.~\hyperref[ch:infotheory]{13}}
\item C.~E.~Shannon, ``Communication theory of secrecy systems,'' \emph{Bell System Technical Journal}, vol.~28, pp.~656--715, 1949.\citedin{Ch.~\hyperref[ch:infotheory]{13}}
\item C.~E.~Shannon, ``Prediction and entropy of printed English,'' \emph{Bell System Technical Journal}, vol.~30, pp.~50--64, 1951.\citedin{Ch.~\hyperref[ch:infotheory]{13}}
\item M.~Stojanovic and J.~Preisig, ``Underwater acoustic communication channels: propagation models and statistical characterization,'' \emph{IEEE Communications Magazine}, vol.~47, no.~1, 2009.\citedin{Ch.~\hyperref[ch:channels]{11}}
\item C.~Sturm and W.~Wiesbeck, ``Waveform design and signal processing aspects for fusion of wireless communications and radar sensing,'' \emph{Proceedings of the IEEE}, vol.~99, no.~7, 2011.\citedin{Ch.~\hyperref[ch:ofdm]{17}}
\item G.~J.~Sullivan \emph{et al.}, ``Overview of the High Efficiency Video Coding (HEVC) standard,'' \emph{IEEE Transactions on Circuits and Systems for Video Technology}, vol.~22, no.~12, 2012.\citedin{Ch.~\hyperref[ch:sourcecoding]{16}}
\item I.~Tal and A.~Vardy, ``List decoding of polar codes,'' \emph{IEEE Transactions on Information Theory}, vol.~61, no.~5, 2015.\citedin{Ch.~\hyperref[ch:moderncodes]{15}}
\item V.~Tarokh, N.~Seshadri and A.~R.~Calderbank, ``Space--time codes for high data rate wireless communication: performance criterion and code construction,'' \emph{IEEE Transactions on Information Theory}, vol.~44, no.~2, 1998.\citedin{Ch.~\hyperref[ch:mimo]{19}}
\item E.~Telatar, ``Capacity of multi-antenna Gaussian channels,'' \emph{European Transactions on Telecommunications}, vol.~10, no.~6, 1999.\citedin{Ch.~\hyperref[ch:mimo]{19}}
\item S.~ten Brink, ``Convergence behavior of iteratively decoded parallel concatenated codes,'' \emph{IEEE Transactions on Communications}, vol.~49, no.~10, 2001.\citedin{Ch.~\hyperref[ch:moderncodes]{15}}
\item F.~A.~Tobagi and L.~Kleinrock, ``Packet switching in radio channels: Part~II,'' \emph{IEEE Transactions on Communications}, 1975.\citedin{Ch.~\hyperref[ch:multipleaccess]{20}}
\item G.~Ungerboeck, ``Channel coding with multilevel/phase signals,'' \emph{IEEE Transactions on Information Theory}, vol.~28, no.~1, 1982.\citedin{Ch.~\hyperref[ch:modulation]{9}}
\item L.~Vangelista, ``Frequency shift chirp modulation: The LoRa modulation,'' \emph{IEEE Signal Processing Letters}, vol.~24, no.~12, 2017.\citedin{Ch.~\hyperref[ch:wifi]{22}}
\item M.~Vanhoef and F.~Piessens, ``Key reinstallation attacks: Forcing nonce reuse in WPA2,'' \emph{Proceedings of ACM CCS}, 2017.\citedin{Ch.~\hyperref[ch:wifi]{22}}
\item S.~Verd\'u, ``Fifty years of Shannon theory,'' \emph{IEEE Transactions on Information Theory}, vol.~44, no.~6, pp.~2057--2078, 1998.\citedin{Ch.~\hyperref[ch:infotheory]{13}}
\item A.~J.~Viterbi, ``Error bounds for convolutional codes and an asymptotically optimum decoding algorithm,'' \emph{IEEE Transactions on Information Theory}, vol.~13, no.~2, 1967.\citedin{Ch.~\hyperref[ch:classiccodes]{14}}
\item J.~E.~Volder, ``The CORDIC trigonometric computing technique,'' \emph{IRE Transactions on Electronic Computers}, vol.~EC-8, no.~3, 1959.\citedin{Ch.~\hyperref[ch:dsp]{6}}
\item R.~H.~Walden, ``Analog-to-digital converter survey and analysis,'' \emph{IEEE Journal on Selected Areas in Communications}, vol.~17, no.~4, 1999.\citedin{Chs.~\hyperref[ch:pcm]{5}, \hyperref[ch:sixg]{25}}
\item G.~K.~Wallace, ``The JPEG still picture compression standard,'' \emph{Communications of the ACM}, vol.~34, no.~4, 1991.\citedin{Ch.~\hyperref[ch:sourcecoding]{16}}
\item Y.-P.~E.~Wang \emph{et al.}, ``A primer on 3GPP narrowband Internet of Things,'' \emph{IEEE Communications Magazine}, vol.~55, no.~3, 2017.\citedin{Ch.~\hyperref[ch:wifi]{22}}
\item D.~K.~Weaver, ``A third method of generation and detection of single-sideband signals,'' \emph{Proceedings of the IRE}, vol.~44, no.~12, 1956.\citedin{Ch.~\hyperref[ch:analog]{4}}
\item S.~B.~Weinstein and P.~M.~Ebert, ``Data transmission by frequency-division multiplexing using the discrete Fourier transform,'' \emph{IEEE Transactions on Communication Technology}, vol.~19, no.~5, 1971.\citedin{Ch.~\hyperref[ch:ofdm]{17}}
\item S.~B.~Weinstein, ``The history of orthogonal frequency-division multiplexing,'' \emph{IEEE Communications Magazine}, vol.~47, no.~11, 2009.\citedin{Ch.~\hyperref[ch:ofdm]{17}}
\item A.~X.~Widmer and P.~A.~Franaszek, ``A DC-balanced, partitioned-block, 8B/10B transmission code,'' \emph{IBM Journal of Research and Development}, vol.~27, no.~5, 1983.\citedin{Ch.~\hyperref[ch:baseband]{8}}
\item T.~Wiegand, G.~J.~Sullivan, G.~Bj{\o}ntegaard and A.~Luthra, ``Overview of the H.264/AVC video coding standard,'' \emph{IEEE Transactions on Circuits and Systems for Video Technology}, vol.~13, no.~7, 2003.\citedin{Ch.~\hyperref[ch:sourcecoding]{16}}
\item P.~J.~Winzer, D.~T.~Neilson and A.~R.~Chraplyvy, ``Fiber-optic transmission and networking: the previous 20 and the next 20 years,'' \emph{Optics Express}, vol.~26, no.~18, 2018.\citedin{Ch.~\hyperref[ch:wireline]{24}}
\item I.~H.~Witten, R.~M.~Neal and J.~G.~Cleary, ``Arithmetic coding for data compression,'' \emph{Communications of the ACM}, vol.~30, no.~6, 1987.\citedin{Ch.~\hyperref[ch:sourcecoding]{16}}
\item Y.~Xiong \emph{et al.}, ``On the fundamental tradeoff of integrated sensing and communications under Gaussian channels,'' \emph{IEEE Transactions on Information Theory}, vol.~69, no.~9, 2023.\citedin{Ch.~\hyperref[ch:sixg]{25}}
\item L.~Zheng and D.~N.~C.~Tse, ``Diversity and multiplexing: a fundamental tradeoff in multiple-antenna channels,'' \emph{IEEE Transactions on Information Theory}, vol.~49, no.~5, 2003.\citedin{Ch.~\hyperref[ch:mimo]{19}}
\end{itemize}
\endgroup

\section*{Standards, specifications and regulations}
\addcontentsline{toc}{section}{Standards, specifications and regulations}
\markright{Standards, specifications and regulations}
\begingroup\small
\begin{itemize}[leftmargin=1.5em,itemindent=-1.5em,itemsep=1.5pt plus 1pt,parsep=0pt,topsep=4pt,label={}]\rightskip=0pt plus 3em\relax
\item 3GPP TR~38.811, \emph{Study on New Radio (NR) to support non-terrestrial networks} (channel models and deployment scenarios).\citedin{Ch.~\hyperref[ch:satellite]{23}}
\item 3GPP TR~38.821, \emph{Solutions for NR to support non-terrestrial networks}.\citedin{Ch.~\hyperref[ch:satellite]{23}}
\item 3GPP TR~38.843, study on artificial intelligence/machine learning for the NR air interface.\citedin{Ch.~\hyperref[ch:sixg]{25}}
\item 3GPP TR~38.901, \emph{Study on channel model for frequencies from 0.5 to 100\,GHz}.\citedin{Chs.~\hyperref[ch:channels]{11}, \hyperref[ch:mimo]{19}}
\item 3GPP TR~45.820, cellular system support for ultra-low complexity and low throughput Internet of Things (the cellular IoT study).\citedin{Ch.~\hyperref[ch:wifi]{22}}
\item 3GPP TS~23.501, \emph{System architecture for the 5G System}.\citedin{Ch.~\hyperref[ch:cellular]{21}}
\item 3GPP TS~25.211--25.214, UTRA (WCDMA) FDD physical layer.\citedin{Ch.~\hyperref[ch:cellular]{21}}
\item 3GPP TS~36.101, E-UTRA user equipment radio transmission and reception (Annex~B: LTE fading propagation conditions).\citedin{Ch.~\hyperref[ch:channels]{11}}
\item 3GPP TS~36.211, E-UTRA physical channels and modulation (Clause~10: NB-IoT).\citedin{Chs.~\hyperref[ch:cellular]{21}, \hyperref[ch:wifi]{22}}
\item 3GPP TS~36.212, E-UTRA multiplexing and channel coding (the LTE turbo code).\citedin{Chs.~\hyperref[ch:moderncodes]{15}, \hyperref[ch:cellular]{21}}
\item 3GPP TS~36.213, E-UTRA physical layer procedures.\citedin{Ch.~\hyperref[ch:cellular]{21}}
\item 3GPP TS~38.101-1, NR user equipment radio transmission and reception.\citedin{Ch.~\hyperref[ch:sdr]{7}}
\item 3GPP TS~38.104, NR base station radio transmission and reception (sensitivity, blocking, ACLR and EVM requirements).\citedin{Chs.~\hyperref[ch:sdr]{7}, \hyperref[ch:modulation]{9}}
\item 3GPP TS~38.211, NR physical channels and modulation (numerology, modulation mapping, PSS/SSS/SSB).\citedin{Chs.~\hyperref[ch:modulation]{9}, \hyperref[ch:sync]{10}, \hyperref[ch:ofdm]{17}, \hyperref[ch:cellular]{21}}
\item 3GPP TS~38.212, \emph{NR; Multiplexing and channel coding}.\citedin{Chs.~\hyperref[ch:moderncodes]{15}, \hyperref[ch:cellular]{21}}
\item 3GPP TS~38.213, NR physical layer procedures for control (SSB, beam failure recovery).\citedin{Chs.~\hyperref[ch:mimo]{19}, \hyperref[ch:cellular]{21}}
\item 3GPP TS~38.214, NR physical layer procedures for data (CSI reporting, Type~I and Type~II codebooks).\citedin{Chs.~\hyperref[ch:mimo]{19}, \hyperref[ch:cellular]{21}}
\item 3GPP TS~38.300, NR overall description (including non-terrestrial networks).\citedin{Ch.~\hyperref[ch:satellite]{23}}
\item 3GPP TS~38.306, NR user equipment radio access capabilities.\citedin{Ch.~\hyperref[ch:cellular]{21}}
\item 3GPP TS~45.002, GSM/EDGE multiplexing and multiple access on the radio path.\citedin{Ch.~\hyperref[ch:cellular]{21}}
\item Bluetooth SIG, \emph{Bluetooth Core Specification}, version~6.0, 2024.\citedin{Ch.~\hyperref[ch:wifi]{22}}
\item CableLabs, CM-SP-PHYv3.1 and CM-SP-PHYv4.0, DOCSIS~3.1 and 4.0 physical layer specifications.\citedin{Ch.~\hyperref[ch:wireline]{24}}
\item CCSDS 131.0-B, \emph{TM Synchronization and Channel Coding} (Blue Book), current issue.\citedin{Ch.~\hyperref[ch:classiccodes]{14}}
\item CEPT ERC Recommendation 70-03, relating to the use of short-range devices (Europe).\citedin{App.~\hyperref[appB]{B}}
\item ETSI EN~300~421, DVB-S: framing structure, channel coding and modulation for 11/12\,GHz satellite services.\citedin{Ch.~\hyperref[ch:classiccodes]{14}}
\item ETSI EN~301~545-2, DVB-RCS2: second generation interactive satellite system.\citedin{Ch.~\hyperref[ch:satellite]{23}}
\item ETSI EN~302~307-1 and -2, DVB-S2 and DVB-S2X: second generation framing structure, channel coding and modulation for satellite services.\citedin{Chs.~\hyperref[ch:modulation]{9}, \hyperref[ch:sync]{10}, \hyperref[ch:classiccodes]{14}, \hyperref[ch:satellite]{23}}
\item IEEE Std 802.3 (Ethernet), current edition, especially Clause~49 (64b/66b PCS), Clause~91 (RS-FEC) and the PAM-4 clauses of 802.3bs, 802.3cd and 802.3ck.\citedin{Chs.~\hyperref[ch:baseband]{8}, \hyperref[ch:wireline]{24}}
\item IEEE Std 802.11-2020, wireless LAN MAC and PHY specifications: Clause~17 (OFDM PHY), Clause~19 (HT), Clause~21 (VHT) and Clause~10 (MAC procedures).\citedin{Ch.~\hyperref[ch:wifi]{22}}
\item IEEE Std 802.11ax-2021, high efficiency WLAN (HE PHY, Clause~27).\citedin{Chs.~\hyperref[ch:modulation]{9}, \hyperref[ch:wifi]{22}}
\item IEEE 802.11be, extremely high throughput WLAN (Wi-Fi~7).\citedin{Chs.~\hyperref[ch:modulation]{9}, \hyperref[ch:wifi]{22}}
\item IEEE Std 802.15.4-2020, \emph{Standard for Low-Rate Wireless Networks}.\citedin{Ch.~\hyperref[ch:wifi]{22}}
\item IEEE Std 1588-2019, Precision Time Protocol (PTP).\citedin{Ch.~\hyperref[ch:sync]{10}}
\item IETF RFC~3550, RTP: a transport protocol for real-time applications.\citedin{Ch.~\hyperref[ch:pcm]{5}}
\item IETF RFC~3551, RTP profile for audio and video conferences (PCMU and PCMA payload types).\citedin{Ch.~\hyperref[ch:pcm]{5}}
\item IS-GPS-200, \emph{NAVSTAR GPS Space Segment/Navigation User Interfaces}, current revision.\citedin{Ch.~\hyperref[ch:spreadspectrum]{18}}
\item IS-GPS-705, GPS L5 interface specification, current revision.\citedin{Ch.~\hyperref[ch:spreadspectrum]{18}}
\item IS-GPS-800, GPS L1C interface specification, current revision.\citedin{Ch.~\hyperref[ch:spreadspectrum]{18}}
\item ITU, \emph{Radio Regulations}, current edition.\citedin{App.~\hyperref[appB]{B}}
\item Recommendation ITU-R M.2083, IMT vision: framework and overall objectives of the future development of IMT for 2020 and beyond (IMT-2020), 2015.\citedin{Ch.~\hyperref[ch:cellular]{21}}
\item Recommendation ITU-R M.2160-0, ``Framework and overall objectives of the future development of IMT for 2030 and beyond,'' ITU, 2023.\citedin{Ch.~\hyperref[ch:sixg]{25}}
\item Report ITU-R M.2410, minimum requirements related to technical performance for IMT-2020 radio interfaces, 2017.\citedin{Chs.~\hyperref[ch:cellular]{21}, \hyperref[ch:sixg]{25}}
\item Recommendation ITU-R P.372, \emph{Radio noise}, current edition.\citedin{Ch.~\hyperref[ch:noise]{3}}
\item Recommendation ITU-R P.453, radio refractive index.\citedin{Ch.~\hyperref[ch:channels]{11}}
\item Recommendation ITU-R P.525, free-space attenuation.\citedin{Ch.~\hyperref[ch:channels]{11}}
\item Recommendation ITU-R P.526, propagation by diffraction.\citedin{Ch.~\hyperref[ch:channels]{11}}
\item Recommendation ITU-R P.530, propagation data and prediction methods for terrestrial line-of-sight links.\citedin{Ch.~\hyperref[ch:channels]{11}}
\item Recommendation ITU-R P.533, HF propagation prediction.\citedin{Ch.~\hyperref[ch:channels]{11}}
\item Recommendation ITU-R P.618, propagation data and prediction methods for Earth--space telecommunication systems.\citedin{Chs.~\hyperref[ch:channels]{11}, \hyperref[ch:satellite]{23}}
\item Recommendation ITU-R P.676, attenuation by atmospheric gases.\citedin{Chs.~\hyperref[ch:channels]{11}, \hyperref[ch:satellite]{23}}
\item Recommendation ITU-R P.837, characteristics of precipitation (rain rate).\citedin{Ch.~\hyperref[ch:satellite]{23}}
\item Recommendation ITU-R P.838, specific attenuation model for rain.\citedin{Chs.~\hyperref[ch:channels]{11}, \hyperref[ch:satellite]{23}}
\item Recommendation ITU-R P.839, rain height model.\citedin{Ch.~\hyperref[ch:satellite]{23}}
\item Recommendation ITU-R P.1238, indoor propagation.\citedin{Ch.~\hyperref[ch:channels]{11}}
\item Recommendation ITU-R P.1411, short-range outdoor propagation.\citedin{Ch.~\hyperref[ch:channels]{11}}
\item Recommendation ITU-R S.465, earth-station antenna reference radiation pattern.\citedin{Ch.~\hyperref[ch:satellite]{23}}
\item Recommendation ITU-R S.580, radiation diagrams for earth-station antennas.\citedin{Ch.~\hyperref[ch:satellite]{23}}
\item ITU-T Recommendation G.168, digital network echo cancellers.\citedin{Ch.~\hyperref[ch:pcm]{5}}
\item ITU-T Recommendation G.652, single-mode optical fibre.\citedin{Ch.~\hyperref[ch:wireline]{24}}
\item ITU-T Recommendation G.694.1, DWDM frequency grid.\citedin{Ch.~\hyperref[ch:wireline]{24}}
\item ITU-T Recommendation G.702, digital hierarchy bit rates.\citedin{Ch.~\hyperref[ch:pcm]{5}}
\item ITU-T Recommendation G.704, frame structures for 1544 and 2048\,kbit/s hierarchical levels (SF, ESF, CRC-4, CAS).\citedin{Ch.~\hyperref[ch:pcm]{5}}
\item ITU-T Recommendation G.707, synchronous digital hierarchy (SDH).\citedin{Ch.~\hyperref[ch:wireline]{24}}
\item ITU-T Recommendation G.709, optical transport network (OTN).\citedin{Ch.~\hyperref[ch:wireline]{24}}
\item ITU-T Recommendation G.711, pulse code modulation of voice frequencies ($\mu$-law and A-law).\citedin{Ch.~\hyperref[ch:pcm]{5}}
\item ITU-T Recommendation G.726, ADPCM.\citedin{Ch.~\hyperref[ch:pcm]{5}}
\item ITU-T Recommendation G.984, gigabit-capable passive optical networks (GPON).\citedin{Ch.~\hyperref[ch:wireline]{24}}
\item ITU-T Recommendation G.992.5, ADSL2+.\citedin{Ch.~\hyperref[ch:wireline]{24}}
\item ITU-T Recommendation G.993.2, VDSL2.\citedin{Ch.~\hyperref[ch:wireline]{24}}
\item ITU-T Recommendation G.993.5, self-FEXT cancellation (vectoring) for VDSL2.\citedin{Ch.~\hyperref[ch:wireline]{24}}
\item ITU-T Recommendation G.8262, synchronous Ethernet equipment clocks (SyncE).\citedin{Ch.~\hyperref[ch:sync]{10}}
\item ITU-T Recommendation G.8271, time and phase synchronization of packet networks.\citedin{Ch.~\hyperref[ch:sync]{10}}
\item ITU-T Recommendation G.8275.1, precision time protocol telecom profile with full timing support.\citedin{Ch.~\hyperref[ch:sync]{10}}
\item ITU-T Recommendation G.9701, G.fast.\citedin{Ch.~\hyperref[ch:wireline]{24}}
\item ITU-T Recommendation G.9807.1, 10-gigabit-capable symmetric passive optical network (\mbox{XGS-PON}).\citedin{Ch.~\hyperref[ch:wireline]{24}}
\item LoRa Alliance, \emph{LoRaWAN L2 1.0.4 Specification}.\citedin{Ch.~\hyperref[ch:wifi]{22}}
\item LoRa Alliance, \emph{RP002 Regional Parameters}.\citedin{Ch.~\hyperref[ch:wifi]{22}}
\item NIST FIPS~203, 204 and 205, post-quantum cryptography standards, 2024.\citedin{Ch.~\hyperref[ch:sixg]{25}}
\item OIF, 400ZR implementation agreement.\citedin{Ch.~\hyperref[ch:wireline]{24}}
\item US Code of Federal Regulations, 47 CFR Part~15 (radio frequency devices: licence-exempt operation).\citedin{App.~\hyperref[appB]{B}}
\item US Code of Federal Regulations, 47 CFR Part~97 (amateur radio service).\citedin{App.~\hyperref[appB]{B}}
\end{itemize}
\endgroup

\section*{Other: application notes, manuals, software and websites}
\addcontentsline{toc}{section}{Other: application notes, manuals, software and websites}
\markright{Other}
\begingroup\small
\begin{itemize}[leftmargin=1.5em,itemindent=-1.5em,itemsep=1.5pt plus 1pt,parsep=0pt,topsep=4pt,label={}]\rightskip=0pt plus 3em\relax
\item Analog Devices, \emph{AD9361/AD9364 Reference Manual} (UG-570), and the AD9364 data sheet.\citedin{Chs.~\hyperref[ch:dsp]{6}, \hyperref[ch:sdr]{7}}
\item Ettus Research, \emph{USRP Hardware Driver and USRP Manual}, \url{https://files.ettus.com/manual/}, and the USRP B200/B210 documentation and product data sheet.\citedin{Ch.~\hyperref[ch:sdr]{7}; App.~\hyperref[appB]{B}}
\item The GNU Radio project, \emph{GNU Radio Wiki and Tutorials}, \url{https://wiki.gnuradio.org}.\citedin{App.~\hyperref[appB]{B}}
\item The jupytext documentation, \url{https://jupytext.readthedocs.io}.\citedin{App.~\hyperref[appB]{B}}
\item Keysight Technologies (formerly Agilent), \emph{Fundamentals of RF and Microwave Noise Figure Measurements}, Application Note~57-1.\citedin{Ch.~\hyperref[ch:noise]{3}}
\item Keysight Technologies (formerly Agilent), \emph{Noise Figure Measurement Accuracy: The Y-Factor Method}, Application Note~57-2.\citedin{Ch.~\hyperref[ch:noise]{3}}
\item Keysight Technologies, \emph{Spectrum Analysis Basics}, Application Note~150 (regularly updated).\citedin{Ch.~\hyperref[ch:signals]{2}}
\item M.~Lichtman, \emph{PySDR: A Guide to SDR and DSP using Python}, \url{https://pysdr.org}.\citedin{App.~\hyperref[appB]{B}}
\item B.~Murmann, ``ADC performance survey,'' online, continuously updated.\citedin{Chs.~\hyperref[ch:pcm]{5}, \hyperref[ch:sixg]{25}}
\item K.~B.~Petersen and M.~S.~Pedersen, \emph{The Matrix Cookbook}, freely available online.\citedin{App.~\hyperref[appA]{A}}
\item \emph{SigMF: the Signal Metadata Format}, \url{https://github.com/sigmf/SigMF}.\citedin{App.~\hyperref[appB]{B}}
\item R.~Volz, \emph{radioconda}, \url{https://github.com/ryanvolz/radioconda}.\citedin{App.~\hyperref[appB]{B}}
\item R.~N.~Williams, ``A painless guide to CRC error detection algorithms,'' 1993.\citedin{Ch.~\hyperref[ch:classiccodes]{14}}
\end{itemize}
\endgroup
